\documentclass[11pt]{article}
\usepackage{amsmath,amssymb,bm}
\usepackage[margin=25mm]{geometry}
\title{V17C.21.1: A Quasi-Steady Adiabatic M3 Slow Flow}
\author{}
\date{}
\begin{document}
\maketitle

\section{State, forcing torus, and harmonic closure}
Let the slow mechanical state of particle $i\in\{1,2\}$ be
\begin{equation}
 X_i=(r_i,u_i,\omega_i,z_i,w_i),\qquad
 u_i=\dot r_i,\quad w_i=\dot z_i,
\end{equation}
and $X=(X_1,X_2)\in\mathbb R^{10}$. The independently driven RF phases are
$\psi_1,\psi_2\in\mathbb T$ with $\dot\psi_1=1$ and
$\dot\psi_2=\nu$ in $\tau=\Omega_1t$. The orbit phase is $\theta_i$.

At a frozen slow state, the retained M3 micromotion coefficients
$\xi_i(\theta_i,\psi_1,\psi_2;X)$ solve the collocation closure
\begin{equation}
 \mathcal G_i(\xi_i;X)=0.
 \label{eq:G}
\end{equation}
The Fourier derivative for index $(n_\theta,n_1,n_2)$ is formed with
$\kappa=n_\theta\omega_i+n_1+n_2\nu$. The fast/slow mask retains RF-dependent
modes with $|\kappa|\ge\kappa_c$ in the fast inverse; the mean and excluded
near-resonant modes remain in the averaged slow residual. The full eight-mode
MFS RF field and the nonlinear Cunningham-corrected Stokes--Schiller--Naumann
drag are evaluated on the phase torus.

\section{Averaged force balance and radial/axial accelerations}
Write $\bar f_{r,i}(X)$ and $\bar f_{z,i}(X)$ for the normalized mean radial
and axial force per unit mass after solving \eqref{eq:G}; these include RF,
DC, drag, gravity, and the ring-averaged Coulomb force. In the rotating frame,
the radial kinematic equation contains centrifugal acceleration. Define
\begin{align}
 a_{r,i}(X)&=r_i\omega_i^2+\bar f_{r,i}(X),\\
 a_{z,i}(X)&=\bar f_{z,i}(X).
\end{align}
The radial and axial components of the slow flow are
\begin{equation}
 \dot r_i=u_i,\quad \dot u_i=a_{r,i}(X),\qquad
 \dot z_i=w_i,\quad \dot w_i=a_{z,i}(X).
 \label{eq:radialaxial}
\end{equation}
The corresponding equilibrium residuals are precisely the radial and axial
M3 force-balance equations used to compute the steady roots.

\section{Angular dynamics from averaged angular momentum}
Let $\mathcal L_i(X)$ be the torus-averaged specific angular momentum returned
by the converged micromotion solution, and let $\mathcal T_i(X)$ be the sum
of mean RF, drag, and Coulomb torques per unit mass. Exact balance of angular
momentum for the reduced description requires
\begin{equation}
 \frac{d\mathcal L_i(X)}{d\tau}=\mathcal T_i(X).
 \label{eq:torque}
\end{equation}
Because $\mathcal L_i$ includes micromotion, it is generally not replaced by
$r_i^2\omega_i$. Applying the chain rule to $X_i=(r_i,u_i,\omega_i,z_i,w_i)$,
\begin{align}
 \dot\omega_i={1\over \partial_{\omega_i}\mathcal L_i}\big[\mathcal T_i
 &-(\partial_{r_i}\mathcal L_i)u_i
 -(\partial_{u_i}\mathcal L_i)a_{r,i}
 -(\partial_{z_i}\mathcal L_i)w_i
 -(\partial_{w_i}\mathcal L_i)a_{z,i}\big].
 \label{eq:omega}
\end{align}
This construction is defined only if $\partial_{\omega_i}\mathcal L_i\ne0$.
On the nominal root, $\partial_{\omega_i}\mathcal L_i$ was positive for both
particles (approximately $0.31305$ and $0.56644$ in the code normalization).
For the ideal infinite-rod field and independent ring-phase Coulomb average,
$\partial_{z_i}\mathcal L_i=0$ at the tested state; it remains explicit in
\eqref{eq:omega} for a nonuniform 3D field.

\section{Implicit dependence of micromotion}
The slow-flow derivative includes how the fast solution changes with $X$.
If $\mathcal G_{i,\xi}$ is invertible on the retained fast subspace and the
HB root is differentiable, then
\begin{equation}
 {\partial\xi_i\over\partial X}
 =-\mathcal G_{i,\xi}^{-1}\mathcal G_{i,X}.
 \label{eq:implicit}
\end{equation}
Rather than form or invert a large dense matrix, this study used central
finite differences of the complete closure: each perturbed slow state triggered
a fresh converged JFNK harmonic solve. Thus the resulting finite-difference
Jacobian includes the numerical response of $\xi_i(X)$, up to inner-solve and
finite-difference errors. The discretized Fourier residual and its convergence
status were checked at every state used in the Jacobian.

\section{Coulomb and averaging assumptions}
The two particles interact through the Coulomb force averaged over independent
uniform ring phases. This yields radial and axial mean forces used in
\eqref{eq:radialaxial}; the phase-averaged tangential Coulomb torque is zero.
It does not describe a phase-locked pair, a close encounter, or a collision.
The RF phases are torus-averaged and therefore are not coordinates of this
autonomous reduced flow.

\section{Local linear stability}
Let $X_*$ be a root of the six circular-orbit balances with $u_i=w_i=0$.
The ten-dimensional slow-flow Jacobian is
\begin{equation}
 J_{M3}=\left.{\partial F_{M3}(X)\over\partial X}\right|_{X_*},
\end{equation}
where $F_{M3}$ is defined by \eqref{eq:radialaxial} and \eqref{eq:omega}.
Eigenvalues are local exponential rates per $\tau$. For the nominal
192$\times$16$\times$16 root, all ten had negative real parts; the largest
was $-0.1059257$. Finite-difference step factors $1/2,1,2$ changed the nearest
matched eigenvalues by less than $1.7\times10^{-6}$ in the reported norm.

\section{Scope and non-claims}
This is a quasi-steady adiabatic closure: the fast equation omits slow
accelerations, the Coulomb term uses independent-phase ring averaging, and the
Fourier fast/slow partition presumes a useful frequency gap. It is not an exact
nonstationary M3 derivation. Orbital angle phases and the two external RF phases
are omitted by averaging, so their neutral directions are not eigenvalues of
$J_{M3}$. Negative eigenvalues establish conditional local stability of this
reduced flow only. They do not prove a full-system attracting torus, global
basin, or Floquet stability.

\end{document}
